%%%%%%%%%%%%%%%%%%%%%%%%%%%%%%%%%%%%%%%%%
% Submission for frontier optics.
% Authors:
% Xiaofeng Qian, Son H. Nguyen, Matthew Daddino, Ching Man Wong
% Stevens Institute of Technology
%%%%%%%%%%%%%%%%%%%%%%%%%%%%%%%%%%%%%%%%%
% !TEX root = Thesis.tex

%----------------------------------------------------------------------------------------
%	PACKAGES AND OTHER DOCUMENT CONFIGURATIONS
%----------------------------------------------------------------------------------------

\documentclass[
	a4paper,
	12pt,
	oneside,
]{article}

\usepackage[utf8]{inputenc}
\usepackage[english]{babel}
\usepackage[]{amsthm}
\usepackage[]{amssymb}
\usepackage[]{setspace}
\usepackage[left=1in, right=1in, top=1in, bottom=1in]{geometry}
\usepackage[colorlinks=true, linkcolor=blue, urlcolor=red, citecolor=blue]{hyperref}
\usepackage{xcolor}
\usepackage{soul}
\usepackage{amsmath}
\usepackage{amsfonts}
\usepackage{amssymb}
\usepackage{graphicx}
\usepackage{float}
\usepackage{braket}
\usepackage{gensymb}
\usepackage{booktabs}

\onehalfspacing

%----------------------------------------------------------------------------------------
%	TITLE AND AUTHOR
%----------------------------------------------------------------------------------------

\title{\textbf{Performing Quantum Computation with Classical Optics: \\[0.5em] A Comprehensive Study of the Variational Quantum Eigensolver Implementation via Photonic Systems}}

\author{Son H. Nguyen\textsuperscript{1}, Matthew Daddino\textsuperscript{1}, Ching Man Wong\textsuperscript{1},Xiaofeng Qian\textsuperscript{1}}

\date{\textsuperscript{1}Physics Department, Stevens Institute of Technology, Hoboken, NJ 07030, USA \\[1em] \today}

%----------------------------------------------------------------------------------------

\begin{document}

\maketitle

%----------------------------------------------------------------------------------------
%	ABSTRACT
%----------------------------------------------------------------------------------------

\begin{abstract}
\noindent Quantum computers promise transformative capabilities for molecular simulation and optimization problems, yet practical implementations face significant technological barriers. This thesis demonstrates an alternative approach: implementing quantum algorithms using classical optical systems. We present a comprehensive framework for performing quantum computation through optical components, specifically realizing the Variational Quantum Eigensolver (VQE) algorithm for molecular energy calculations. Our implementation encodes qubits using photonic polarization states and spatial modes, leveraging wave plates (quarter-wave and half-wave), beamsplitters, and dove prisms to construct quantum gates. We detail the theoretical decomposition of the UCCSD (Unitary Coupled Cluster Single and Double excitations) ansatz for the hydrogen molecule, deriving the necessary rotations and entangling operations. The optical realization employs Jones calculus to systematically map quantum gates onto polarization-dependent optical transformations, with Stokes vector measurements enabling quantum state readout. This comprehensive study bridges quantum-algorithm design with optical laboratory realization, demonstrating that quantum algorithms can be executed through photonic platforms as an alternative to superconducting or trapped-ion systems. All derivations, circuit decompositions, and optical implementations are provided in full mathematical detail.
\end{abstract}
% Minimal ending to allow standalone compilation
\clearpage
\begin{thebibliography}{99}
\bibitem{omalley2016scalable} O'Malley, P. J. J., et al. (2016). Scalable quantum simulation of molecular energies. \textit{Physical Review X}, 6(3), 031007.
\end{thebibliography}

\end{document}